\begin{lemma}[The signed power of a partial derivative for $1<p<2$, and inequality (3.12)]
Let $n\ge2$, $1<p<2$ and $r>0$.  Then there is a constant
$C=C(n,p,r)>0$ with the following property.  Let
$\varepsilon>0$, $\gamma_0>0$, $\gamma_1\le r\gamma_0$, let
$x_0\in\R^n$ and $R>0$, write $B_\rho=B_\rho(x_0)$, let $\mu>0$ with
\[
\varepsilon<\Bigl(\frac\mu2\Bigr)^{2},
\]
let $A$ satisfy (Q1) of the thesis--(Q2) of the thesis on $B_R$ with parameters
$p,\varepsilon,\gamma_0,\gamma_1$ (\noderef{EllipticStructureQ12}), assume the
cited regularity result (C1) for $\varepsilon>0$
(\noderef{CitedRegularityEpsPos}), let $u$ be a weak solution of
$\operatorname{div}A(\nabla u)=0$ in $B_R$ (\noderef{IsWeakSolutionOn}) with weak
gradient $g=\nabla u$, suppose that for each $i$ the component $g^i=D_iu$
belongs to $W^{1,2}_{\mathrm{loc}}(B_R)$ with weak gradient $Dg^i$
(\noderef{MemW1pLocOn}), and suppose
\[
|g^i|\le\mu\quad\text{almost everywhere on }B_R,\qquad i=1,\dots,n .
\]
Fix an index $k$ and put
\[
\lambda=\frac\mu2,\qquad h_0=\lambda^{p/2},\qquad
v=|g^k|^{\frac p2-1}g^k,\qquad
V=\Bigl(\frac p2|g^k|^{\frac p2-1}\Bigr)Dg^k .
\]
Then both of the following hold.
\begin{enumerate}
\item For every $T$ with $0<T<R$ one has $v\in W^{1,2}(B_T)$ with weak gradient
  $V$ (\noderef{MemW1pOn}).
\item For all $\rho$ with $0<\rho\le R$, all $\sigma$ with $0<\sigma<1$ and all
  $h\in\R$:
  \[
  h\ge-h_0\ \Longrightarrow\
  \int_{B_{\rho-\sigma\rho}\cap\{v>h\}}\|V\|^2\,dx
  \ \le\ C\,\frac{h_0^2}{(\sigma\rho)^2}\,\bigl|B_\rho\cap\{v>h\}\bigr| ,
  \]
  \[
  h\le h_0\ \Longrightarrow\
  \int_{B_{\rho-\sigma\rho}\cap\{v<h\}}\|V\|^2\,dx
  \ \le\ C\,\frac{h_0^2}{(\sigma\rho)^2}\,\bigl|B_\rho\cap\{v<h\}\bigr| .
  \]
\end{enumerate}

\medskip
\noindent\textbf{Why the exponent restriction is a case and not a convenience.}
The same two conclusions hold for $p\ge2$, and the reader may ask why they are
not stated once for every $p>1$.  The answer is that the subquadratic range is
not a special case of a uniform argument but a \emph{different} argument, and
the three things that differ are the three things an argument is made of.
\emph{The auxiliary function differs.}  For $p\ge2$ the map
$s\mapsto|s|^{p/2-1}s$ is itself $C^1$ --- its derivative $\tfrac
p2|s|^{p/2-1}$ is continuous exactly because $p/2-1\ge0$ --- so $v$ is a $C^1$
function of $g^k$ and nothing else is needed.  For $1<p<2$ that derivative blows
up at $0$, the chain rule is unavailable, and the thesis introduces the
regularised pair $\pi_\eta(s)=(|s|+\eta)^{p-2}s$ and
$G_\eta(s)=\operatorname{sgn}(s)[(|s|+\eta)^{p/2}-\eta^{p/2}]$ together with a
second auxiliary function $\psi=|g^k|^{p-2}g^k$, which has no counterpart in the
other range.  \emph{The source inequality differs.}  For $p\ge2$ conclusion (2)
is read off the weighted Caccioppoli inequality (3.4) at a level of $g^k$.  For
$1<p<2$ that inequality is the wrong shape --- the weight enters with a negative
power and the truncation is of $\psi$, not of $g^k$ --- and what is used instead
is the inequality (3.17) proved below, at a level of $\psi$.  \emph{The producer
node differs.}  (3.4) comes from the node \texttt{CaccioppoliLevelWeight},
which this lemma therefore never uses; (3.17) is proved here from
\noderef{DifferentiatedEquationWeak} directly, against a test
function built from the regularised truncation, and passes to the limit
$\eta\downarrow0$ by Fatou and dominated convergence.  Nothing is shared by the
two ranges except the hypothesis package and the elementary transport of levels,
and in particular no weakening of the constant $C$ lets one argument cover both.
That is why this range is a node of its own rather than a branch inside one.

\medskip
The hypothesis $\varepsilon<(\mu/2)^2$ is the degenerate regime, in which the
weight $w=\varepsilon+|g|^2$ is pinched between $\lambda^2$-sized constants on
the level sets rather than between $\varepsilon$-sized ones; both conclusions
are consumed only in that regime.  The minimality of $\mu$ is not assumed: only
the bound $|g^i|\le\mu$ is used.  The function $v$ is the thesis'
$|g^k|^{p/2}\operatorname{sgn}(g^k)$ written without a sign function: the two
agree at every real value of $g^k$, the factor $g^k$ killing the value of the
real power $|0|^{p/2-1}$ at a zero of $g^k$.  Likewise $V$ is written as a
single real power, and on $\{g^k=0\}$ the run's convention gives
$|0|^{p/2-1}=0$, which for $1<p<2$ is the value the thesis prescribes there.
\end{lemma}
\begin{proof}
Throughout, $w=\varepsilon+|g|^2$, and $C_1,C_2,\dots$ denote positive
constants depending only on $n$, $p$ and $r$; this is checked at each
occurrence.  Note $\lambda=\mu/2>0$ and $h_0=\lambda^{p/2}>0$.

\emph{Step A: four preliminaries.}

\textbf{(A1) The two-sided bound on $w$.}  Almost everywhere on $B_R$ we have
$|g^i|\le\mu=2\lambda$ for every $i$, hence
$|g|^2=\sum_i(g^i)^2\le 4n\lambda^2$, and therefore
\[
\varepsilon\ \le\ w\ \le\ \varepsilon+4n\lambda^2\ <\ (1+4n)\lambda^2
\qquad\text{almost everywhere on }B_R,
\tag{A.1}
\]
the last inequality by $\varepsilon<\lambda^2$.  Both bounds are used only under
integral signs, so the almost-everywhere form is the form needed.

\textbf{(A2) $\gamma_0\le\gamma_1$.}  The clauses \texttt{Q1} and \texttt{Q2} of
\noderef{EllipticStructureQ12} each hold for almost every $x\in B_R$, hence
simultaneously at almost every such $x$, and $|B_R|>0$ (Preamble: the volume of
a ball, $R>0$); fix such an $x$.  Take $z$ arbitrary, $h=0$ and $\xi=e_1$ the
first standard basis vector --- available because $n\ge2\ge1$ --- noting
$\|e_1\|=1$; the guard
$\varepsilon+\|0\|^2=\varepsilon\ne0$ holds because $\varepsilon>0$.  Then
\texttt{Q1}, whose left-hand side is
$\gamma_0\varepsilon^{p/2-1}\|e_1\|^2=\gamma_0\varepsilon^{p/2-1}$, gives
$\gamma_0\varepsilon^{p/2-1}\le\sum_i\bigl(DA(0)e_1\bigr)_i(e_1)_i
=\bigl(DA(0)e_1\bigr)_1$, the sum collapsing because $(e_1)_i$ vanishes for
$i\ne1$ and is $1$ for $i=1$; and \texttt{Q2}, at $j=i=1$, gives
$\bigl|\bigl(DA(0)e_1\bigr)_1\bigr|\le\gamma_1\varepsilon^{p/2-1}$, so
$\gamma_0\varepsilon^{p/2-1}\le\gamma_1\varepsilon^{p/2-1}$; dividing by
$\varepsilon^{p/2-1}>0$ gives $\gamma_0\le\gamma_1$.  This is the form in which
\noderef{DifferentiatedEquationWeak} asks for the two constants, while the
present lemma assumes $\gamma_1\le r\gamma_0$.

\textbf{(A3) Localisation.}  Let $0<T<R$.  Then $\operatorname{closure}B_T
=\overline B_T$ is compact and contained in $B_R$, so clause (4) of
\noderef{WeakGradientRestrict}, applied with $q=2$ and the open sets $B_R$ and
$B_T$, upgrades $g^k\in W^{1,2}_{\mathrm{loc}}(B_R)$ with weak gradient $Dg^k$
(\noderef{MemW1pLocOn}) to
\[
g^k\in W^{1,2}(B_T)\ \text{with weak gradient }Dg^k\qquad(0<T<R).
\tag{A.2}
\]
Conversely, if a pair $(f,G)$ satisfies $f\in W^{1,2}(B_T)$ with weak gradient
$G$ for every $T<R$, then $f\in W^{1,2}_{\mathrm{loc}}(B_R)$ with weak gradient
$G$: the defining identity of \noderef{HasWeakGradientOn} on $B_R$ quantifies
over test functions $\varphi$ with $\operatorname{tsupport}\varphi\subseteq B_R$
compact, and every such $\operatorname{tsupport}\varphi$ lies in some $B_T$ with
$T<R$, where the identity is available and where the integrals over $B_T$ and
over $B_R$ agree because both integrands vanish off
$\operatorname{tsupport}\varphi$; and every compact $K\subseteq B_R$ lies in such
a $B_T$, so the $L^2$ bounds required by \noderef{MemW1pLocOn} follow from those
supplied by \noderef{MemW1pOn} on $B_T$.  We refer to this passage as (A.2$'$).
Finally $\operatorname{volume}B_T\ne\infty$ for every $T$ (Preamble: the volume
of a ball), which is the finiteness hypothesis of
\noderef{WeakGradientC1Comp}, \noderef{WeakGradientPositivePart} and clause (1)
of \noderef{WeakGradientSmoothMul}.

\textbf{(A4) The cut-off.}  Let $\rho_{1/8}$ be the standard mollifier of the
Preamble in dimension one at scale $\tfrac18$ --- smooth, nonnegative, supported
in $[-\tfrac18,\tfrac18]$, of total integral $1$ --- and put
$\psi_\flat=\mathbf1_{(-\infty,1/2]}*\rho_{1/8}$.  Then
$\psi_\flat\in C^\infty(\R)$ with
$\psi_\flat'=\mathbf1_{(-\infty,1/2]}*\rho_{1/8}'$ (Preamble: convolution with a
smooth compactly supported kernel, the indicator being locally integrable);
$0\le\psi_\flat\le1$, because $0\le\mathbf1_{(-\infty,1/2]}\le1$ and
$\rho_{1/8}\ge0$ has total integral $1$, so the convolution is an average of
values in $[0,1]$; $\psi_\flat=1$ on $(-\infty,\tfrac38]$ and $\psi_\flat=0$ on
$[\tfrac58,\infty)$, since writing the convolution as
$\int\mathbf1_{(-\infty,1/2]}(t-s)\rho_{1/8}(s)\,ds$ (a change of variable
legitimate by the translation and reflection invariance of Lebesgue measure,
Preamble) one finds for $t\le\tfrac38$ that every $s$ with $|s|\le\tfrac18$ has
$t-s\le\tfrac12$, so the integrand is $\rho_{1/8}$ on its whole support and the
integral is $1$, while for $t\ge\tfrac58$ every such $s$ has $t-s\ge\tfrac12$,
so the indicator vanishes except possibly at the single point $s=t-\tfrac12$, a
Lebesgue-null set, and the integral is $0$; and
$|\psi_\flat'|\le\|\rho_{1/8}'\|_{L^1(\R)}=:\Lambda<\infty$, an absolute
constant.  Given $0<\rho_0<\rho_1$ put
\[
\zeta_{\rho_0,\rho_1}(x)=\psi_\flat
\Bigl(\frac{|x-x_0|-\rho_0}{\rho_1-\rho_0}\Bigr).
\]
Then $0\le\zeta\le1$; $\zeta=1$ on $\overline B_{\rho_0}$; $\zeta=0$ outside
$\overline B_{\rho_0+\frac58(\rho_1-\rho_0)}$, so $\operatorname{supp}\zeta$ is
compact and contained in $B_{\rho_1}$; $\zeta$ is smooth, because
$x\mapsto|x-x_0|$ is smooth away from $x_0$ and $\zeta$ is constant on the
neighbourhood $B_{\rho_0}$ of $x_0$; and
$|\nabla\zeta|\le\Lambda/(\rho_1-\rho_0)$, with $\nabla\zeta$ supported in
$B_{\rho_1}$.

\emph{Step B: the level correspondence.}  The map
$s\mapsto|s|^{p/2-1}s$ is a strictly increasing bijection of $\R$ onto $\R$: it
is odd, it vanishes only at $0$, and on $(0,\infty)$ it is $s\mapsto s^{p/2}$
with $p/2>0$, which is a strictly increasing bijection onto $(0,\infty)$.  Hence
for every $\ell\in\R$, with
$h=|\ell|^{p/2-1}\ell$,
\[
\{g^k>\ell\}=\{v>h\},\qquad \{g^k<\ell\}=\{v<h\},
\tag{B.1}
\]
and $|\ell|\le\lambda\iff|h|\le h_0$, because $|h|=|\ell|^{p/2}$ and
$h_0=\lambda^{p/2}$ and $t\mapsto t^{p/2}$ is increasing on $[0,\infty)$.
Moreover $|v|=|g^k|^{p/2}\le(2\lambda)^{p/2}=2^{p/2}h_0$ almost everywhere on
$B_R$, by the hypothesis $|g^k|\le\mu=2\lambda$.

\emph{Step D: conclusion (1) and inequality (3.17) for $1<p<2$.}  Here
$s\mapsto|s|^{p/2-1}s$ is not $C^1$ at $0$ and the thesis' regularisation is
needed.  For $\eta>0$ put
\[
\pi_\eta(s)=(|s|+\eta)^{p-2}s,\qquad
G_\eta(s)=\operatorname{sgn}(s)\bigl[(|s|+\eta)^{p/2}-\eta^{p/2}\bigr],
\qquad \psi_\eta=\pi_\eta\circ g^k,\qquad v_\eta=G_\eta\circ g^k .
\]
Both $\pi_\eta$ and $G_\eta$ are odd and strictly increasing, and both are $C^1$
with bounded derivative:
\[
\pi_\eta'(s)=(|s|+\eta)^{p-2}+(p-2)(|s|+\eta)^{p-3}|s|,
\qquad
G_\eta'(s)=\tfrac p2(|s|+\eta)^{p/2-1},
\]
the first continuous at $0$ because its second summand carries the factor $|s|$,
and both bounded by $\eta^{p-2}$ resp.\ $\tfrac p2\eta^{p/2-1}$, since for
$1<p<2$ the exponents $p-2$ and $p/2-1$ are negative and $|s|+\eta\ge\eta$.  One
also has
\[
(p-1)(|s|+\eta)^{p-2}\ \le\ \pi_\eta'(s)\ \le\ (|s|+\eta)^{p-2} ,
\tag{D.1}
\]
the upper bound because $(p-2)(|s|+\eta)^{p-3}|s|\le0$, the lower because
$(|s|+\eta)^{p-3}|s|\le(|s|+\eta)^{p-2}$ and hence
$\pi_\eta'(s)\ge(|s|+\eta)^{p-2}+(p-2)(|s|+\eta)^{p-2}=(p-1)(|s|+\eta)^{p-2}$.
These are the thesis' inequalities on p.~21.  Fix $0<T<R$.  By (A.2), (A3) and
\noderef{WeakGradientC1Comp} on $B_T$ with $q=2$, --- that node's composition
hypothesis being met because $\pi_\eta$ and $G_\eta$ are globally $C^1$ with
bounded derivative, so that no almost-everywhere transport of the pair is
needed ---
\[
\psi_\eta\in W^{1,2}(B_T)\ \text{with weak gradient }\pi_\eta'(g^k)Dg^k,
\qquad
v_\eta\in W^{1,2}(B_T)\ \text{with weak gradient }
V_\eta:=\tfrac p2(|g^k|+\eta)^{p/2-1}Dg^k ,
\tag{D.2}
\]
so that $\|V_\eta\|^2=\tfrac{p^2}4(|g^k|+\eta)^{p-2}\|Dg^k\|^2$, which is the
thesis' (3.14) of the thesis.  Note also
$|\psi_\eta|\le(|g^k|+\eta)^{p-2}|g^k|\le|g^k|^{p-1}\le(2\lambda)^{p-1}$ almost
everywhere, using $p-2<0$, and $|v_\eta|\le(2\lambda+\eta)^{p/2}$.

\textbf{(D3) The differentiated equation against the regularised truncation.}
Fix $\ell\in\R$, $0<\sigma<1$ and $0<\rho\le R$, and let
$\zeta=\zeta_{\rho-\sigma\rho,\rho}$ of (A4).  Put
\[
\theta^+_\eta=(\psi_\eta-\ell)^+,\qquad \varphi^+=\theta^+_\eta\,\zeta^2 .
\]
Let $T<R$ be any radius with $\operatorname{tsupport}\zeta\subseteq B_T$; such a
$T$ exists because $\operatorname{tsupport}\zeta$ is a compact subset of
$B_\rho\subseteq B_R$.  By \noderef{WeakGradientPositivePart} with $q=2$, the
level $c=\ell$, the open set $B_T$ of finite volume and the pair
$(\psi_\eta,\pi_\eta'(g^k)Dg^k)$ of (D.2), $\theta^+_\eta\in W^{1,2}(B_T)$ with
weak gradient $\chi_{\{\psi_\eta>\ell\}}\pi_\eta'(g^k)Dg^k$.  Since (D.2) holds
for every $T<R$ and the construction of $\theta^+_\eta$ does not involve
$\zeta$, (A.2$'$) gives $\theta^+_\eta\in W^{1,2}_{\mathrm{loc}}(B_R)$ with that
weak gradient.  Now $\zeta^2$ is smooth with compact support in $B_R$, so clause
(2) of \noderef{WeakGradientSmoothMul} with $q=2$ and $\Omega=B_R$ gives
$\varphi^+\in W^{1,2}(B_R)$ with weak gradient
\[
g_{\varphi^+}=\zeta^2\chi_{\{\psi_\eta>\ell\}}\pi_\eta'(g^k)Dg^k
+\theta^+_\eta\,\nabla(\zeta^2),
\qquad \nabla(\zeta^2)=2\zeta\nabla\zeta ,
\]
and $\varphi^+$ has compact support with
$\operatorname{tsupport}\varphi^+\subseteq\operatorname{tsupport}\zeta\subseteq
B_R$.  Apply \noderef{DifferentiatedEquationWeak} with
$n,p,\varepsilon,\gamma_0,\gamma_1$ --- its hypotheses $2\le n$, $1<p$,
$0<\varepsilon$, $0<\gamma_0$ and $\gamma_0\le\gamma_1$, the last by (A2) ---
the open set $B_R$, the field $A$, the solution $u$ with gradient $g$, the
family $Dg$, the index $k$ and the pair $(\varphi^+,g_{\varphi^+})$.  It gives
\[
\int_{B_R}\sum_{i,j}\bigl(DA(g)e_j\bigr)_i\,(Dg^k)_j\,(g_{\varphi^+})_i\,dx=0 ,
\]
Separating the two summands of $g_{\varphi^+}$, and using that for each fixed
$i$ one has $\sum_j\bigl(DA(g)e_j\bigr)_i(Dg^k)_j=\bigl(DA(g)Dg^k\bigr)_i$,
because $DA(g)$ is linear in its direction argument and
$Dg^k=\sum_j(Dg^k)_je_j$, this reads
\[
\int_{B_R}\chi_{\{\psi_\eta>\ell\}}\pi_\eta'(g^k)\,
\bigl\langle DA(g)Dg^k,Dg^k\bigr\rangle\zeta^2
=-2\int_{B_R}\theta^+_\eta\,\zeta\,
\bigl\langle DA(g)Dg^k,\nabla\zeta\bigr\rangle .
\]
On the left, \texttt{Q1} of \noderef{EllipticStructureQ12} with $\xi=Dg^k$ and
$h=g(x)$ --- its guard $\varepsilon+\|g\|^2=w\ne0$ holding since
$\varepsilon>0$ --- gives
$\langle DA(g)Dg^k,Dg^k\rangle\ge\gamma_0w^{p/2-1}\|Dg^k\|^2$, and (D.1) gives
$\pi_\eta'(g^k)\ge(p-1)(|g^k|+\eta)^{p-2}$.  On the right, \texttt{Q2} gives
$|(DA(g)e_j)_i|\le\gamma_1w^{p/2-1}$, whence
\[
\bigl|\langle DA(g)Dg^k,\nabla\zeta\rangle\bigr|
\le\gamma_1w^{\frac p2-1}\sum_{i,j}|(Dg^k)_j|\,|(\nabla\zeta)_i|
\le n\gamma_1w^{\frac p2-1}\|Dg^k\|\,|\nabla\zeta| ,
\]
using $\sum_j|a_j|\le\sqrt n\,\|a\|$ (Cauchy--Schwarz, Preamble) twice.  Hence
\[
\gamma_0(p-1)\!\!\int_{A^+_\ell(\eta)}\!\!
w^{\frac p2-1}\|Dg^k\|^2(|g^k|+\eta)^{p-2}\zeta^2
\ \le\ 2n\gamma_1\!\!\int_{B_R}\!\!
w^{\frac p2-1}\|Dg^k\|\,\theta^+_\eta\,|\zeta|\,|\nabla\zeta| ,
\tag{D.3}
\]
where $A^+_\ell(\eta)=B_R\cap\{\psi_\eta>\ell\}$; the integrand on the right
vanishes off $A^+_\ell(\eta)$ because $\theta^+_\eta$ does.  This is the thesis'
(3.13) of the thesis.

\textbf{(D4) The Cauchy--Schwarz step.}  Assume in addition
$\eta^2<\varepsilon$.  Then $w\ge\varepsilon>\eta^2$ and $w\ge(g^k)^2$, so
$w^{1/2}\ge\max\{\eta,|g^k|\}\ge\tfrac12(|g^k|+\eta)$ and therefore, since
$1-p/2>0$,
\[
(|g^k|+\eta)^{1-\frac p2}\le\bigl(2w^{1/2}\bigr)^{1-\frac p2}
=2^{1-\frac p2}w^{\frac{2-p}4}.
\]
From (D.2), $\|Dg^k\|=\tfrac2p(|g^k|+\eta)^{1-\frac p2}\|V_\eta\|$, so
\[
w^{\frac p2-1}\|Dg^k\|\ \le\ \tfrac2p2^{1-\frac p2}\,
w^{\frac p2-1+\frac{2-p}4}\|V_\eta\|\ =\ C_1\,w^{\frac{p-2}4}\|V_\eta\| ,
\tag{D.4}
\]
with $C_1=C_1(p)=\tfrac2p2^{1-p/2}$, because
$\tfrac p2-1+\tfrac{2-p}4=\tfrac{p-2}4$.  Substituting (D.4) and the identity
$\|V_\eta\|^2=\tfrac{p^2}4(|g^k|+\eta)^{p-2}\|Dg^k\|^2$ into (D.3), and applying
the Cauchy--Schwarz inequality (Preamble) to the right-hand side written as the
product of $w^{\frac{p-2}4}\|V_\eta\|\,|\zeta|$ and
$\theta^+_\eta|\nabla\zeta|$ --- and noting
$\bigl(w^{\frac{p-2}4}\bigr)^2=w^{\frac p2-1}$ ---
\[
\tfrac{4\gamma_0(p-1)}{p^2}\!\!\int_{A^+_\ell(\eta)}\!\!
w^{\frac p2-1}\|V_\eta\|^2\zeta^2
\ \le\ 2n\gamma_1C_1\Bigl(\int_{A^+_\ell(\eta)}\!\!
w^{\frac p2-1}\|V_\eta\|^2\zeta^2\Bigr)^{\frac12}
\Bigl(\int_{B_R}(\theta^+_\eta)^2|\nabla\zeta|^2\Bigr)^{\frac12}.
\]
The first factor on the left is finite: $\|V_\eta\|\le\tfrac
p2\eta^{p/2-1}\|Dg^k\|$ and $\|Dg^k\|$ is in $L^2$ of the compact set
$\operatorname{tsupport}\zeta$ by \noderef{MemW1pLocOn}, while
$w^{p/2-1}\le\varepsilon^{p/2-1}$ since $p/2-1<0$ and $w\ge\varepsilon$, and
$\zeta$ is bounded and vanishes off $\operatorname{tsupport}\zeta$.  So that
factor may be divided out; squaring, and using $\gamma_1\le r\gamma_0$ to
cancel $\gamma_0$,
\[
\int_{A^+_\ell(\eta)}w^{\frac p2-1}\|V_\eta\|^2\zeta^2\,dx
\ \le\ C_2\int_{B_R}\bigl((\psi_\eta-\ell)^+\bigr)^2|\nabla\zeta|^2\,dx ,
\tag{D.5}
\]
with $C_2=\bigl(2nrC_1p^2/(4(p-1))\bigr)^2$, a function of $n,p,r$.  The $-$
case, with $(\ell-\psi_\eta)^+$ and $A^-_\ell(\eta)=B_R\cap\{\psi_\eta<\ell\}$,
follows by the same argument from $\varphi^-=(\ell-\psi_\eta)^+\zeta^2$: by
\noderef{WeakGradientLinear} (clause (2), $q=2$, $a=-1$, $b=0$, both pairs taken
to be the pair of (D.2)) the pair $(-\psi_\eta,-\pi_\eta'(g^k)Dg^k)$ lies in
$W^{1,2}(B_T)$, and \noderef{WeakGradientPositivePart} applied to it at the
level $-\ell$ gives $(\ell-\psi_\eta)^+\in W^{1,2}(B_T)$ with weak gradient
$-\chi_{\{\psi_\eta<\ell\}}\pi_\eta'(g^k)Dg^k$; the remaining estimates are
identical, the sign of the test function's gradient entering only through the
absolute values taken above.

\textbf{(D5) The limit $\eta\downarrow0$: conclusion (1).}  Put
\[
\psi=|g^k|^{p-2}g^k\ \text{on}\ \{g^k\ne0\},\quad\psi=0\ \text{on}\ \{g^k=0\} .
\]
We first claim $Dg^k=0$ almost everywhere on $\{g^k=0\}$.  Indeed on $B_T$,
\noderef{WeakGradientPositivePart} at the level $c=0$ gives $(g^k)^+\in
W^{1,2}(B_T)$ with weak gradient $\chi_{\{g^k>0\}}Dg^k$, and applied to
$(-g^k,-Dg^k)$ --- which lies in $W^{1,2}(B_T)$ by clause (2) of
\noderef{WeakGradientLinear} with $a=-1$, $b=0$ --- it gives $(-g^k)^+$ with
weak gradient $-\chi_{\{g^k<0\}}Dg^k$; since $g^k=(g^k)^+-(-g^k)^+$, clause (2)
of \noderef{WeakGradientLinear} with $a=1$, $b=-1$ shows that
$\chi_{\{g^k>0\}}Dg^k+\chi_{\{g^k<0\}}Dg^k$ is also a weak gradient of $g^k$ on
$B_T$; and a weak gradient is determined almost everywhere (Preamble: the
fundamental lemma of the calculus of variations, second consequence).  Hence
$\chi_{\{g^k=0\}}Dg^k=0$ almost everywhere, as claimed.  Consequently
$\|V_\eta\|\uparrow\|V\|$ pointwise almost everywhere as $\eta\downarrow0$: on
$\{g^k\ne0\}$ because $\eta\mapsto(|g^k|+\eta)^{p/2-1}$ increases to
$|g^k|^{p/2-1}$ as $\eta\downarrow0$, the exponent being negative, and on
$\{g^k=0\}$ because both sides are $0$ almost everywhere there.

Take $\ell=0$ in (D.5) and its $-$ partner: then $A^+_0(\eta)=B_R\cap\{g^k>0\}$
and $A^-_0(\eta)=B_R\cap\{g^k<0\}$ for every $\eta$, because $\pi_\eta$ is odd
and strictly increasing so $\{\psi_\eta>0\}=\{g^k>0\}$; and
$\bigl((\psi_\eta)^\pm\bigr)^2\le(2\lambda)^{2(p-1)}$ almost everywhere.  Fix
$T<R$ and $\zeta=\zeta_{T,T'}$ with $T<T'<R$; then $\zeta=1$ on $\overline B_T$
and $|\nabla\zeta|\le\Lambda/(T'-T)$, so (D.5), its $-$ partner and
$w^{p/2-1}\ge\bigl((1+4n)\lambda^2\bigr)^{p/2-1}>0$ --- valid almost everywhere
by (A.1) and $p/2-1<0$ --- give
\[
\int_{B_T\cap\{g^k\ne0\}}\|V_\eta\|^2
\ \le\ \bigl((1+4n)\lambda^2\bigr)^{1-\frac p2}C_2
\frac{\Lambda^2}{(T'-T)^2}(2\lambda)^{2(p-1)}\,2|B_R|\ =:\ \mathcal K<\infty
\]
for every $\eta$ with $\eta^2<\varepsilon$.  By Fatou's lemma (Preamble) and
$\|V_\eta\|\uparrow\|V\|$ we get $\int_{B_T}\|V\|^2\le\mathcal K<\infty$.
Next, restrict to $\eta\le1$: then $v_\eta\to v$ pointwise, because
$G_\eta(s)\to|s|^{p/2}\operatorname{sgn}(s)=|s|^{p/2-1}s$ for every $s$, with
$|v_\eta|\le(2\lambda+1)^{p/2}$ almost everywhere; and $V_\eta\to V$ pointwise
almost everywhere with $\|V_\eta\|\le\|V\|\in L^2(B_T)$.  Hence for every test
function $\varphi$ with $\operatorname{tsupport}\varphi\subseteq B_T$ compact,
dominated convergence (Preamble) --- the dominating functions being
$(2\lambda+1)^{p/2}|D_i\varphi|$ and $\|V\|\,|\varphi|$, both integrable on the
finite-measure set $B_T$ --- passes to the limit in the identity of
\noderef{HasWeakGradientOn} for $(v_\eta,V_\eta)$ and yields it for $(v,V)$.
Together with $|v|\le2^{p/2}h_0$ almost everywhere (Step B) and
$\|V\|\in L^2(B_T)$ this is $v\in W^{1,2}(B_T)$ with weak gradient $V$
(\noderef{MemW1pOn}), which is conclusion (1) for $1<p<2$.

\textbf{(D6) The limit $\eta\downarrow0$: inequality (3.17).}  Write
$A^+_\ell=B_R\cap\{\psi>\ell\}$.  Fatou's lemma needs only a pointwise lower
bound on the liminf of the integrands, and \emph{not} any monotonicity of the
sets $A^+_\ell(\eta)$; what has to be checked is the pointwise statement
\[
x\in A^+_\ell\ \Longrightarrow\ x\in A^+_\ell(\eta)\ \text{ for all sufficiently
small }\eta>0 ,
\tag{D.6}
\]
and it holds for \emph{every} real $\ell$.  Let $x\in B_R$ with $\psi(x)>\ell$
and put $t=g^k(x)$, so $\psi(x)=|t|^{p-2}t$ and $\psi_\eta(x)=(|t|+\eta)^{p-2}t$;
recall $p-2<0$.  Three cases.
\begin{itemize}
\item $t>0$: then $\eta\mapsto(t+\eta)^{p-2}t$ is continuous at $\eta=0$,
  because $t>0$ keeps the base away from $0$, with value $\psi(x)>\ell$; so
  $\psi_\eta(x)>\ell$ for all small $\eta>0$.
\item $t=0$: then $\psi(x)=0$ and $\psi_\eta(x)=\eta^{p-2}\cdot0=0$ for every
  $\eta$, the factor $t$ vanishing whatever value the real power $\eta^{p-2}$
  takes.  So $\psi(x)>\ell$ says $\ell<0$ and $\psi_\eta(x)=0>\ell$ for every
  $\eta>0$.
\item $t<0$: then $|t|+\eta>|t|>0$ and $p-2<0$ give
  $(|t|+\eta)^{p-2}\le|t|^{p-2}$; multiplying by $t<0$ reverses the inequality,
  so $\psi_\eta(x)\ge\psi(x)>\ell$ for every $\eta>0$.
\end{itemize}
It is worth recording what is \emph{not} true, since it is what a
monotonicity-based argument would need: on $\{g^k>0\}$ the same computation run
at $t>0$ gives $\psi_\eta\le\psi$, so the sets $A^+_\ell(\eta)$ do not increase
to $A^+_\ell$.  That costs nothing, because Fatou consumes (D.6) and not
inclusion.

Fix $\ell\in\R$, $0<\sigma<1$, $0<\rho\le R$ and $\zeta=\zeta_{\rho-\sigma\rho,\rho}$,
and apply Fatou's lemma (Preamble) to the nonnegative integrands
$\chi_{A^+_\ell(\eta_j)}w^{p/2-1}\|V_{\eta_j}\|^2\zeta^2$ along any sequence
$\eta_j\downarrow0$ with $\eta_j^2<\varepsilon$: at each $x\in A^+_\ell$ the
characteristic factor equals $1$ for all large $j$ by (D.6), and
$\|V_{\eta_j}(x)\|\uparrow\|V(x)\|$ almost everywhere by (D5), so the liminf of
the integrands is at least $\chi_{A^+_\ell}w^{p/2-1}\|V\|^2\zeta^2$ almost
everywhere; hence the integral of the latter is at most the liminf of the
left-hand sides of (D.5).  Meanwhile $(\psi_{\eta}-\ell)^+\to(\psi-\ell)^+$
pointwise: at a point with $t\ne0$ the map $\eta\mapsto(|t|+\eta)^{p-2}t$ is
continuous at $\eta=0$ with value $\psi(x)$, and at a point with $t=0$ one has
$\psi_\eta(x)=0=\psi(x)$ for every $\eta$, so $\psi_\eta\to\psi$ pointwise and
$s\mapsto(s-\ell)^+$ is continuous.  The bound
$\bigl((\psi_\eta-\ell)^+\bigr)^2\le\bigl((2\lambda)^{p-1}+|\ell|\bigr)^2$ holds
almost everywhere, by $|\psi_\eta|\le|g^k|^{p-1}\le(2\lambda)^{p-1}$, and
$|\nabla\zeta|^2$ is bounded and supported in $B_\rho$, hence integrable; so
dominated convergence (Preamble) identifies the limit of the right-hand sides of
(D.5).  Therefore, for every $\ell\in\R$,
\[
\int_{A^+_\ell}w^{\frac p2-1}\|V\|^2\zeta^2\,dx
\ \le\ C_2\int_{B_R}\bigl((\psi-\ell)^+\bigr)^2|\nabla\zeta|^2\,dx ,
\tag{3.17$^+$}
\]
and symmetrically (3.17$^-$) for every $\ell\in\R$, with
$A^-_\ell=B_R\cap\{\psi<\ell\}$ and $(\ell-\psi)^+$ in place of
$(\psi-\ell)^+$; there the three cases of (D.6) are $t<0$ (continuity, since
$|t|>0$), $t=0$ (then $\psi(x)=0<\ell$ forces $\ell>0$ and $\psi_\eta(x)=0<\ell$
for every $\eta$) and $t>0$ (where $\psi_\eta\le\psi<\ell$ for every $\eta$).

\emph{Step E: conclusion (2).}  Fix $0<\rho\le R$ and $0<\sigma<1$ and let
$\zeta=\zeta_{\rho-\sigma\rho,\rho}$ of (A4), so that
$|\nabla\zeta|\le\Lambda/(\sigma\rho)$, $\zeta=1$ on
$\overline B_{\rho-\sigma\rho}$, $0\le\zeta\le1$ and
$\operatorname{supp}\zeta\subseteq B_\rho$.  In each case below the level $h$ is
transported to a level $\ell$ by a strictly increasing bijection, and only a
one-sided bound on $\ell$ is used.

\textbf{(E3) Case $1<p<2$, the $-$ side.}  Let $h\le h_0$ and put
\[
\ell=|h|^{\frac{2(p-1)}p-1}h ,
\qquad\text{equivalently}\qquad
h=|\ell|^{\frac p{2(p-1)}-1}\ell .
\]
Both maps are strictly increasing bijections of $\R$ (odd, vanishing only at
$0$, and on $(0,\infty)$ equal to $t\mapsto t^{2(p-1)/p}$ resp.\
$t\mapsto t^{p/(2(p-1))}$, which are mutually inverse increasing bijections of
$(0,\infty)$ since the two exponents are reciprocal).  Composing,
$\{\psi<\ell\}=\{v<h\}$: indeed $\psi=|g^k|^{p-2}g^k$ and $v=|g^k|^{p/2-1}g^k$
are both strictly increasing functions of $g^k$, and $\ell$ and $h$ are the
images of the same point of $\R$ under the corresponding bijections, namely
$\psi<\ell\iff g^k<|\ell|^{\frac1{p-1}-1}\ell\iff v<h$.  The level transports
exactly:
\[
h\le h_0=\lambda^{p/2}\iff\ell\le\lambda^{p-1},
\qquad\text{because}\qquad
\bigl(\lambda^{p-1}\bigr)^{\frac p{2(p-1)}}=\lambda^{\frac p2}
\]
and $t\mapsto|t|^{\frac p{2(p-1)}-1}t$ is increasing.  Now use (3.17$^-$) at
this $\ell$, available for \emph{every} real $\ell$ by (D6).  Its left-hand
side: by (A.1), $\varepsilon<\lambda^2$ and $p/2-1<0$ we have
$w^{p/2-1}\ge\bigl((1+4n)\lambda^2\bigr)^{p/2-1}=(1+4n)^{p/2-1}\lambda^{p-2}$
almost everywhere, and $\zeta=1$ on $\overline B_{\rho-\sigma\rho}$, so
\[
(1+4n)^{\frac p2-1}\lambda^{p-2}
\int_{B_{\rho-\sigma\rho}\cap\{v<h\}}\|V\|^2
\ \le\ \int_{A^-_\ell}w^{\frac p2-1}\|V\|^2\zeta^2 .
\]
Its right-hand side: the integrand is supported in $\{\psi<\ell\}\cap B_\rho$,
where $0<\ell-\psi\le\ell^++(2\lambda)^{p-1}\le\lambda^{p-1}
+2^{p-1}\lambda^{p-1}=(1+2^{p-1})\lambda^{p-1}$, using $\ell\le\lambda^{p-1}$
and $|\psi|\le(2\lambda)^{p-1}$ almost everywhere, so
\[
C_2\int_{B_R}\bigl((\ell-\psi)^+\bigr)^2|\nabla\zeta|^2
\le C_2(1+2^{p-1})^2\lambda^{2p-2}\frac{\Lambda^2}{(\sigma\rho)^2}
\bigl|B_\rho\cap\{v<h\}\bigr| .
\]
Dividing both sides by $(1+4n)^{p/2-1}\lambda^{p-2}>0$ --- strictly positive for
every real $p$ because $\lambda>0$, including the present case where $p-2<0$ ---
and using $\lambda^{2p-2}/\lambda^{p-2}=\lambda^p=h_0^2$ gives the second
inequality of conclusion (2) with
$C=C_2(1+2^{p-1})^2(1+4n)^{1-p/2}\Lambda^2$.

\textbf{(E4) Case $1<p<2$, the $+$ side.}  Let $h\ge-h_0$ and put the same
$\ell$, so that $\{\psi>\ell\}=\{v>h\}$ and $\ell\ge-\lambda^{p-1}$ by the
mirror image of the transport identity.  Run the same computation from
(3.17$^+$) at this $\ell$ --- again available for every real $\ell$ --- the
right-hand bound now being
$0<\psi-\ell\le(2\lambda)^{p-1}+\ell^-\le(1+2^{p-1})\lambda^{p-1}$ with
$\ell^-=\max\{-\ell,0\}\le\lambda^{p-1}$.  This gives the first inequality of
conclusion (2).

Taking for $C$ the maximum of the two expressions produced in (E3) and (E4) ---
each a function of $n$, $p$ and $r$ alone, since $C_2$ and $\Lambda$ are
--- both inequalities of conclusion (2) hold on the full one-sided ranges
$h\ge-h_0$ and $h\le h_0$ with one constant.  Conclusion (1) was proved in
(D5).  This completes the proof.
\end{proof}
